\chapter{Functions, Devices and Agents}
\label{ch:faas-iot-agentic}

The web and LLM tasks of \Cref{ch:web-tasks,ch:llm-tasks} call services that
answer one request with one reply. This chapter covers three task types that
reach further (Lectures~29--36). \texttt{FAAS} tasks invoke cloud functions
that are already deployed. \texttt{IOT} tasks observe and act on physical
devices. \texttt{AGENTIC} tasks run a supervised agent against a declared
contract. All three share a design principle that sets them apart from batch
tasks: their checkpoints are judged by \emph{content}, not by name, because
their effects are paid for, physical, or non-deterministic.

\section{FaaS Tasks}

\subsection{The Contract}

A \texttt{FAAS} task names a provider (\texttt{mock}, \texttt{http},
\texttt{knative}, \texttt{aws}, \texttt{azure}, \texttt{gcp}, or one added
with \texttt{register\_provider}) and a deployed function. It declares its
inputs, outputs, invocation mode, timeout and retry policy. DAGonStar does not
deploy functions. It invokes them through provider adapters, and cloud SDKs
are imported only when their provider is used.

\lstinputlisting[caption={Lecture 29: a FAAS task against the mock provider (\texttt{examples/29-faas/faas\_mock.py}).},label={lst:faas-mock}]{../../examples/29-faas/faas_mock.py}

\lstinputlisting[style=bash,caption={Recorded output of Lecture 29.}]{../../examples/29-faas/output/faas_mock.txt}

\subsection{Envelope, Descriptors and Idempotency}

Each invocation sends a CloudEvents-style envelope (\texttt{specversion 1.0},
type \texttt{org.dagonstar.task.invoke}). Its \texttt{id} is the SHA-256 of
the workflow name, the task name and the specification. The id is the same on
every retry and every resume, so a function can recognise a repeated request.
File inputs are not passed as paths. They become \emph{descriptors} with the
original reference, the staged path, the size, the media type, a checksum and
a transport. Retries apply only to transient error codes (\texttt{throttled},
\texttt{provider\_unavailable}, \texttt{transport\_error}). A function that
raises is a permanent failure.

\begin{definitionbox}{Exactly-once results}
A completed FAAS task records the SHA-256 of its specification. On re-execution it is reused only if the recorded exit code is 0, the
specification hash is unchanged and the folder exists. A changed specification is invoked again, and an unchanged one never is.
Lecture~29 proves this by replacing the provider factory with a counting mock
and observing zero calls.
\end{definitionbox}

Lecture~30 runs the same contract against an HTTP function on
\texttt{127.0.0.1} (\texttt{examples/30-faas-http/faas\_http.py}). Only the
provider and its options change.

\section{IoT Tasks}

\subsection{Bounded Operations}

An \texttt{IOT} task performs one of four operations: \texttt{observe},
\texttt{configure}, \texttt{actuate} or \texttt{compute}. It must declare a
\emph{bounded} completion mode: \texttt{first\_event},
\texttt{sample\_count}, \texttt{duration}, \texttt{condition},
\texttt{acknowledged}, \texttt{reported\_state} or \texttt{job}. Observation
modes that wait for something need \texttt{timeout\_seconds}. A task with no
mode is rejected at construction with ``IoT completion must select a
supported bounded mode''.

\lstinputlisting[caption={Lecture 31: a bounded observation (\texttt{examples/31-iot-observe/iot\_observe.py}).},label={lst:iot-observe}]{../../examples/31-iot-observe/iot_observe.py}

\lstinputlisting[style=bash,caption={Recorded output of Lecture 31.}]{../../examples/31-iot-observe/output/iot_observe.txt}

Five readings are offered and \texttt{sample\_count} keeps three, whose
mean is $(16.0 + 16.5 + 17.0)/3 = 16.5$\,\textdegree C. Each task writes
\texttt{.dagon/iot\_provenance.json} with the requested and selected target,
the placement, the migration history, protocol metadata and the outcome
certainty. An observation is reused only while every output file still has
its recorded SHA-256. Deleting the output makes the checkpoint unusable,
because an observation cannot be repeated.

\subsection{Actuation at Most Once}

\lstinputlisting[caption={Lecture 32: configuration and actuation (\texttt{examples/32-iot-actuate/iot\_actuate.py}).},label={lst:iot-actuate}]{../../examples/32-iot-actuate/iot_actuate.py}

\lstinputlisting[style=bash,caption={Recorded output of Lecture 32.}]{../../examples/32-iot-actuate/output/iot_actuate.txt}

\begin{cautionbox}[Unknown outcomes are never replayed]
An actuation is delivered at most once by default. \texttt{max\_attempts}
greater than one is refused unless the delivery declares
\texttt{idempotent: true}. When an actuation is not acknowledged, its
checkpoint records \texttt{outcome\_certainty: unknown}, and the task fails
with ``automatic replay is disabled''. A resume then reattaches to the
recorded invocation, if the provider supports it, rather than sending a
second command. A person decides what happened to the harbour gate.
\end{cautionbox}

\section{Agentic Tasks}

\subsection{Profiles and Contracts}

An agentic task never contains an endpoint, a model, a key or code. It names
a \emph{profile}, a \texttt{[agent.<name>]} section of \texttt{dagon.ini} that
sets the provider, the strategy (\texttt{bounded\_iterative}), the permitted
tools and guardrails, and finite budgets (\texttt{max\_steps},
\texttt{timeout\_seconds}, \texttt{max\_tool\_calls}, \texttt{max\_tokens},
\texttt{max\_attempts}). The task may only narrow tools, lower budgets and
add guardrails. It declares a \emph{contract}: a specification with
objective, context and output schema; inputs as \texttt{workflow://}
references; requirements scored by registered evaluators; and an acceptance
threshold.

\lstinputlisting[caption={The course's quality-control contract (\texttt{examples/common/agentic.py}, excerpt).},label={lst:agentic-common},firstline=76,lastline=93]{../../examples/common/agentic.py}

\lstinputlisting[style=bash,caption={Recorded output of Lecture 33 (\texttt{examples/33-agentic-contract/agentic\_contract.py}).}]{../../examples/33-agentic-contract/output/agentic_contract.txt}

With no tools, the deterministic mock adapter first submits the table
unchanged. The \texttt{record\_coverage} evaluator scores it $4/\max(4,5) =
0.8$, below the threshold of 1.0. After this feedback the agent submits a
deduplicated table, which scores 1.0 and is published. Four invalid
configurations are refused before any adapter, tool or transfer starts.

\subsection{Guardrails}

Every proposed action passes runtime invariants before it executes. The tool
must be permitted and registered. The arguments must be an object that
matches the tool's schema. The arguments may contain no \texttt{workflow://}
reference and no absolute or parent-directory path. Named guardrails
(\texttt{preserve\_record\_ids}, \texttt{non\_empty\_result},
\texttt{read\_only\_tools}) check inputs, actions and results. A repairable
violation becomes feedback for the agent; any other violation fails the task
with \texttt{guardrail\_blocked}.

\lstinputlisting[style=bash,caption={Recorded output of Lecture 34 (\texttt{examples/34-agentic-guardrails/agentic\_guardrails.py}).}]{../../examples/34-agentic-guardrails/output/agentic_guardrails.txt}

\subsection{Evaluation, Publication and Identity}

The result fit is a weighted mean of requirement scores:
\[
F = \frac{\sum_{i:\,w_i>0} w_i s_i}{\sum_{i:\,w_i>0} w_i}.
\]
A candidate is accepted only if no evaluator failed, every hard check passed,
$F$ reaches the threshold, and no mandatory requirement scored below its
minimum. A result is published only after acceptance, through a commit
manifest written last. A checkpoint is reused only if its \emph{identity
digest} still matches. The digest covers the contract, the effective
profile, the tool, guardrail, evaluator and adapter versions, and the SHA-256
of every input.

\lstinputlisting[caption={Lecture 35: gates and content-based reuse (\texttt{examples/35-agentic-evaluation/agentic\_evaluation.py}).},label={lst:agentic-eval}]{../../examples/35-agentic-evaluation/agentic_evaluation.py}

\lstinputlisting[style=bash,caption={Recorded output of Lecture 35.}]{../../examples/35-agentic-evaluation/output/agentic_evaluation.txt}

The golden case is worth computing by hand: $F = (2\cdot1 + 1\cdot0)/3 =
0.667$ exceeds the threshold of 0.6, yet the candidate is rejected because the
mandatory \texttt{stations} requirement failed.

\section{One Graph Across the Continuum}

Lecture~36 joins an IoT observation at a gateway, a native normalisation, the
agentic quality check and a FaaS archive into one chain. References work the
same way in every task type, and each task is checkpointed by its own type's
rule.

\lstinputlisting[caption={Lecture 36: edge to cloud (\texttt{examples/36-continuum/continuum.py}).},label={lst:continuum}]{../../examples/36-continuum/continuum.py}

\lstinputlisting[style=bash,caption={Recorded output of Lecture 36.}]{../../examples/36-continuum/output/continuum.txt}

\begin{table}[h]
\centering
\begin{tabular}{@{}lll@{}}
\toprule
Task type & Reused when & Why \\
\midrule
BATCH, NATIVE & code 0 and folder exists & cheap to repeat \\
FAAS & same specification SHA-256 & paid calls \\
IOT & every output SHA-256 matches & observations are unrepeatable \\
AGENTIC & identity digest and outputs verify & non-deterministic \\
\bottomrule
\end{tabular}
\caption{Checkpoint reuse rules by task type.}
\label{tab:reuse-rules}
\end{table}

\section{Chapter Summary}

\begin{itemize}[leftmargin=1.4em]
  \item FAAS tasks invoke deployed functions with a stable idempotency key, file descriptors and exactly-once results.
  \item IOT tasks are always bounded. Actuations are sent at most once, and
    unknown outcomes are recorded, never replayed.
  \item AGENTIC tasks run inside a profile's authority, are checked before
    every action, and publish only accepted results.
  \item Content-based checkpoints protect effects that are expensive,
    physical or non-deterministic.
\end{itemize}

\begin{exercise}
Extend the continuum workflow of Lecture~36 with an actuation that sounds a
siren when the accepted table contains a temperature above a threshold. Decide
its delivery policy and justify it.
\end{exercise}
